\chapter{Templates}\label{chap:format}

\section{Main Body}

The main body of a thesis is usually organized into chapters, sections, subsections, and subsubsections. The English mode uses English fixed labels and chapter names while keeping the same page geometry and thesis structure.

\section{Figures}

Figures are placed in the \texttt{Img} directory and referenced by file name. Figure~\ref{fig:english-sample} shows a simple example.

\begin{figure}[!htbp]
    \centering
    \includegraphics[width=0.40\textwidth]{tc_q_criteria}
    \bicaption{Q判据等值面图。}{Isocontour of Q criteria.}
    \label{fig:english-sample}
\end{figure}

\section{Tables}

Table~\ref{tab:english-sample} gives a minimal table example.

\begin{table}[!htbp]
    \bicaption{样表示例。}{A sample table.}
    \centering
    \footnotesize
    \begin{tabular}{lccc}
        \hline
        Item & A & B & C \\
        \hline
        Row 1 & 1 & 2 & 3 \\
        Row 2 & 4 & 5 & 6 \\
        \hline
    \end{tabular}
    \label{tab:english-sample}
\end{table}

\section{Equations}

Equations are numbered by chapter. For example,

\begin{equation}
    E = mc^2 .
\end{equation}

\section{Algorithms}

Algorithm~\ref{alg:english-euclid} demonstrates the algorithm environment.

\begin{algorithm}[!htbp]
    \small
    \caption{Euclid's algorithm}\label{alg:english-euclid}
    \begin{algorithmic}[1]
        \Procedure{Euclid}{$a,b$}
        \While{$b \ne 0$}
            \State $r \gets a \bmod b$
            \State $a \gets b$
            \State $b \gets r$
        \EndWhile
        \State \Return $a$
        \EndProcedure
    \end{algorithmic}
\end{algorithm}

\section{References}

Citations are written with commands such as \verb|\cite{lamport1986document}|. For example, \LaTeX{} is introduced in \cite{lamport1986document}.
